\begin{exercises}
    \subsection{Bài tập}

    \begin{questions}[series=SGK]
        \item Cho tam giác $ABC$ vuông tại $A$. Tính các tỉ số lượng giác sin, côsin, tang, côtang của các góc nhọn $B$ và $C$ khi biết:
        \begin{questions}
            \item $AB=8$ cm, $BC=17$ cm;
            \item $AC=0{,}9$ cm, $AB=1{,}2$ cm.
        \end{questions}

        \answerline
        \answerline

        \item Cho tam giác vuông có một góc nhọn $60^\circ$ và cạnh kề với góc $60^\circ$ bằng $3$ cm. Hãy tính cạnh đối của góc này.

        \answerline

        \item Cho tam giác vuông có một góc nhọn bằng $30^\circ$ và cạnh đối với góc này bằng $5$ cm. Tính độ dài cạnh huyền của tam giác.

        \answerline

        \item Cho hình chữ nhật có chiều dài và chiều rộng lần lượt là $3$ và $\sqrt3$. Tính góc giữa đường chéo và cạnh ngắn hơn của hình chữ nhật (sử dụng bảng giá trị lượng giác trang 69).

        \answerline

        \item
        \begin{questions}
            \item Viết các tỉ số lượng giác sau thành tỉ số lượng giác của các góc nhỏ hơn $45^\circ$:
            \[
                \sin55^\circ,\quad \cos62^\circ,\quad \tan57^\circ,\quad \cot64^\circ.
            \]

            \answerline
            \item Tính $\dfrac{\tan25^\circ}{\cot65^\circ}$, $\tan34^\circ-\cot56^\circ$.

            \answerline
        \end{questions}

        \item Dùng MTCT, tính (làm tròn đến chữ số thập phân thứ ba):
        \begin{questions}
            \item $\sin40^\circ12'$;
            \item $\cos52^\circ54'$;
            \item $\tan63^\circ36'$;
            \item $\cot25^\circ18'$.
        \end{questions}

        \answerline

        \item Dùng MTCT, tìm số đo của góc nhọn $x$ (làm tròn đến phút), biết rằng:
        \begin{questions}
            \item $\sin x=0{,}2368$;
            \item $\cos x=0{,}6224$;
            \item $\tan x=1{,}236$;
            \item $\cot x=2{,}154$.
        \end{questions}

        \answerline
    \end{questions}
\end{exercises}
